\clearpage
\phantomsection
\addcontentsline{toc}{chapter}{Cronograma Semestral Integrado (16 Semanas — 2026/2)}
\chapter*{Cronograma Semestral e Plano de Estudos — 16 Semanas}
\markboth{CRONOGRAMA — 16 SEMANAS}{CRONOGRAMA — 16 SEMANAS}

O componente curricular \textbf{CHF662 — Análise Ambiental} totaliza \textbf{60 horas-aula} distribuídas ao longo de \textbf{16 semanas letivas} no período 2026/2 (encontros semanais de 4 horas, das 07h30 às 11h30, com suporte de atendimento no Laboratório GEOTRÓPICOS).

O quadro a seguir integra integralmente o conteúdo programático oficial do Plano de Disciplina à ordem dos capítulos desta apostila, indicando as leituras obrigatórias, as entregas dos Estudos Orientados (EO1–EO4), a Avaliação Teórica, a Viagem de Campo e a construção progressiva do Trabalho Prático Final.

\begin{boxConceito}[Rotina Recomendada de Estudo Autônomo Semanal (4h em sala + 3h extraclasse)]
\begin{itemize}[leftmargin=1.5em, itemsep=2.5pt]
  \item \textbf{Antes da aula (60 min):} leia as Etapas 1 a 3 do capítulo da semana e anote dúvidas sobre os conceitos e símbolos novos.
  \item \textbf{Durante o encontro presencial (4h):} participe da exposição dialogada, estudos dirigidos, oficinas cartográficas e discussão dos textos-base.
  \item \textbf{Após a aula (90 min):} refaça os Estudos de Caso (Etapa 4), resolva os Exercícios de Fixação (Etapa 6) e confira o Gabarito Comentado (Etapa 8).
  \item \textbf{Revisão e Projeto (30 min):} complete o Teste de Aprendizagem (Etapa 10), registre eventuais erros conceituais no Diário de Revisão e alimente o banco de dados do seu Diagnóstico Ambiental Aplicado.
\end{itemize}
\end{boxConceito}

\section*{Cronograma Oficial Semana a Semana (CHF662 — 2026/2)}

\renewcommand{\arraystretch}{1.22}
\small
\begin{longtable}{@{}c p{2.9cm} p{5.5cm} p{4.4cm}@{}}
  \toprule
  \rowcolor{uefsNavy}
  \textcolor{white}{\bfseries Sem.} & \textcolor{white}{\bfseries Módulo / Cap.} & \textcolor{white}{\bfseries Conteúdo Programático e Atividade} & \textcolor{white}{\bfseries Referências-Base e Entregas} \\
  \midrule
  \endfirsthead
  \toprule
  \rowcolor{uefsNavy}
  \textcolor{white}{\bfseries Sem.} & \textcolor{white}{\bfseries Módulo / Cap.} & \textcolor{white}{\bfseries Conteúdo Programático e Atividade} & \textcolor{white}{\bfseries Referências-Base e Entregas} \\
  \midrule
  \endhead
  \bottomrule
  \endfoot

  \textbf{1} & \textbf{Mód. 1 — Cap. 1} \newline Questão ambiental & Apresentação da disciplina e do plano. Introdução à questão ambiental contemporânea, crise ambiental, solos, água, uso da terra e papel da Geografia. & Leff (2001, cap. 2); Porto-Gonçalves (2006). \\
  \rowcolor{bgGray}
  \textbf{2} & \textbf{Mód. 1 — Cap. 2} \newline Sociedade–natureza & A modernidade, a Geografia e a relação sociedade–natureza; natureza como construção social; conflitos socioambientais. & Porto-Gonçalves (2006, p. 8–36); Mendonça (2001). \textbf{\color{uefsAlert}Entrega: EO1.} \\
  \textbf{3} & \textbf{Mód. 1 — Cap. 3} \newline Abordagens geogr. & Abordagens ambientais no pensamento geográfico moderno; Geografia Socioambiental, ambiente, território, paisagem e interdisciplinaridade. & Mendonça (2001); Bertrand (2004). \textbf{\color{uefsAlert}Entrega: EO2.} \\
  \rowcolor{bgGray}
  \textbf{4} & \textbf{Mód. 1 — Cap. 4} \newline Conceitos de AA & Ambiente, impacto ambiental, degradação, vulnerabilidade, fragilidade ambiental, risco, resiliência, recuperação e sustentabilidade. & Sánchez (2013, cap. 1); Tauk-Tornisielo (1996). \textbf{\color{uefsAlert}Entrega: EO3 + Simulado M1.} \\
  \midrule
  \textbf{5} & \textbf{Mód. 2 — Cap. 5} \newline AIA e planejamento & Avaliação de Impacto Ambiental: diagnóstico, prognóstico, identificação de impactos, mitigação, compensação e monitoramento. & Sánchez (2013, cap. 7); Brasil/MMA (2016). \textbf{\color{uefsAlert}Entrega: EO4.} \\
  \rowcolor{bgGray}
  \textbf{6} & \textbf{Mód. 2 — Cap. 6} \newline Análise integrada & Elementos físico, biótico e antrópico; escalas espaciais e temporais; sistemas ambientais; planejamento e ordenamento territorial. & Castro; Lemos (2016); Arruda (2013); Drew (1989). \textbf{Simulado M2.} \\
  \midrule
  \textbf{7} & \textbf{Mód. 3 — Cap. 7} \newline Bacia hidrográfica & Microbacias, divisores, rede de drenagem, morfometria, áreas ripárias, conectividade hidrossedimentológica, erosão e uso da terra. & Botelho; Silva (2004); Christofoletti (1981); Guerra; Cunha (1996). \\
  \rowcolor{bgGray}
  \textbf{8} & \textbf{Mód. 3 — Cap. 8} \newline Morfopedologia & Processos interativos relevo–solo–água: declividade, cobertura vegetal, drenagem, erosão, movimentos de massa e assoreamento. & Guerra (2018); Casseti (2007); Ross (2006). \\
  \textbf{9} & \textbf{Mód. 3 — Cap. 9} \newline Geossist. e Ross & Geossistema, compartimentação geoambiental, ecodinâmica (estável, intergrade, instável), fragilidade potencial e emergente. & Bertalanffy (2010); Bertrand (2004); Tricart (1977); Ross (1994). \\
  \rowcolor{bgTeal}
  \textbf{10} & \textbf{Avaliação Teórica} \newline Simulado Mód. 3 & \textbf{Prova escrita dissertativa presencial (3,0 pontos):} abrangendo integralmente os Capítulos 1 a 9 (Módulos 1, 2 e 3). & Síntese: Leff, Porto-Gonçalves, Mendonça, Sánchez, Bertrand, Tricart e Ross. \\
  \midrule
  \textbf{11} & \textbf{Mód. 4 — Cap. 10} \newline Campo e Indicadores & \textbf{Viagem de Campo (Semiárido / Canudos e eixo energético):} matriz eólica na Caatinga, supressão vegetal, erosão na base das torres e avifauna; indicadores PEIR. & Fidalgo (2003); IBGE (2018); OECD (2003); Roteiro de Campo DCHF/UEFS. \\
  \rowcolor{bgGray}
  \textbf{12} & \textbf{Mód. 4 — Cap. 11} \newline Qualidade do solo & Indicadores físicos, químicos e biológicos do solo: MO, CO, densidade, porosidade, infiltração, agregação, resistência à penetração e funcionalidade edáfica. & Doran; Parkin (1994); Karlen et al. (1997); Vezzani; Mielniczuk (2009). \\
  \textbf{13} & \textbf{Mód. 4 — Cap. 12} \newline Degradação & Diagnóstico de degradação: erosão laminar, sulcos, ravinas, voçorocas, compactação, instabilidade de margens, assoreamento e degradação ripária. & Guerra; Cunha (1996); Ross (1994); Cavalcanti (2006). \textbf{Simulado M4.} \\
  \midrule
  \rowcolor{bgGray}
  \textbf{14} & \textbf{Mód. 5 — Cap. 13} \newline Bioengenharia & Princípios da bioengenharia de água e solo; técnicas vegetativas e estruturais leves; paliçadas, biomantas, geotêxteis, retentores de sedimentos e revegetação. & Gray; Sotir (1996); Schiechtl; Stern (1996); Durlo; Sutili (2005). \\
  \textbf{15} & \textbf{Mód. 5 — Cap. 14} \newline Legislação e PRAD & PNMA (Lei 6.938/81), Código Florestal (Lei 12.651/12), APPs, licenciamento ambiental, PRAD, recomposição vegetal e responsabilidade ambiental. & Brasil (1981; 2012); Brasil/MMA (2016); Sánchez (2013). \textbf{Simulado M5.} \\
  \rowcolor{bgTeal}
  \textbf{16} & \textbf{Mód. 5 — Cap. 15} \newline Encerramento & \textbf{Entrega e Socialização do Trabalho Prático Final (5,0 pontos):} diagnósticos ambientais, matriz de indicadores e propostas de intervenção. & Bibliografia específica de cada diagnóstico + Simulado Geral Final. \\
\end{longtable}
\normalsize

\section*{Diário de Revisão Conceitual e Acompanhamento do Diagnóstico}

Para consolidar a aprendizagem ao longo das 16 semanas — especialmente na preparação para a Avaliação Teórica (Semana 10) e para a defesa do Trabalho Prático Final (Semana 16) — utilize o quadro de registro abaixo:

\begin{center}
\renewcommand{\arraystretch}{1.55}
\small
\begin{tabularx}{\textwidth}{|c|p{2.6cm}|p{3.0cm}|X|c|}
  \hline
  \rowcolor{uefsNavy}
  \textcolor{white}{\bfseries Semana} & \textcolor{white}{\bfseries Cap. / Tema} & \textcolor{white}{\bfseries Dúvida ou Erro} & \textcolor{white}{\bfseries Correção Conceitual / Aplicação no Trabalho Final} & \textcolor{white}{\bfseries Dominado?} \\
  \hline
  \textit{Ex.: 4} & \textit{Cap. 4 — Risco} & \textit{Confundi perigo com risco} & \textit{Perigo é a ameaça física; risco exige exposição de população/bens vulneráveis ($R = P \times V$).} & \textit{[ X ]} \\
  \hline
  & & & & [ \quad ] \\
  \hline
  & & & & [ \quad ] \\
  \hline
  & & & & [ \quad ] \\
  \hline
\end{tabularx}
\end{center}
